\begin{table}[H]
    \centering
    \caption{Representative data and metadata types in catalytic experiments.}
    \label{tab:catalytic-data-landscape}
    \small
    \renewcommand{\arraystretch}{1.2}
    \begin{tabularx}{\textwidth}{p{0.23\textwidth}p{0.35\textwidth}X}
        \toprule
        Experimental area & Typical data and metadata & Role in interpretation \\
        \midrule
        Catalyst material & Composition or formula, metal loading, support material, and unique sample identifier & Identifies the material under study and links results to a specific catalyst sample \cite{schumann2025nomad,salim2026methane}. \\
        Synthesis and treatment & Preparation method, precursor identity, calcination or reduction parameters, and treatment atmosphere & Preserves the catalyst history, since preparation and activation steps define the material state used in testing \cite{mendes2025data,salim2026methane}. \\
        Physicochemical characterization & Surface area, pore properties, particle or crystallite size, and surface composition & Connects structural and chemical descriptors to catalytic behavior and material comparison \cite{schumann2025nomad,salim2026methane}. \\
        Reactor and setup & Reactor type, reactor dimensions, catalyst bed mass, diluent, and reactor filling & Provides the experimental context needed to interpret kinetic measurements and compare setups \cite{wulf2021unified,schumann2025nomad}. \\
        Operating conditions & Reaction temperature and pressure, feed composition, total flow rate, and space velocity or contact time & Defines the process conditions under which performance values were measured \cite{schumann2025nomad,salim2026methane}. \\
        Performance and stability & Conversion, selectivity, yield, reaction rate, and time-on-stream behavior & Reports the functional outcome of the experiment and supports comparison across catalysts and conditions \cite{schumann2025nomad,salim2026methane}. \\
        \bottomrule
    \end{tabularx}
\end{table}
